\subsection{Task 9: Tỉ trọng cơ cấu loại chủ nhà}

\subsubsection{Thiết kế Idiom}

\begin{table}[H]
\centering
\renewcommand{\arraystretch}{1.1}
\begin{tabularx}{\textwidth}{|l|l|X|}
\hline
\textbf{Idiom} & \multicolumn{2}{l|}{\textbf{Pie Chart (Biểu đồ tròn)}} \\ \hline
\textbf{What}
& \multicolumn{2}{X|}{Loại chủ nhà (Host Type): Categorical (C)} \\ \cline{2-3}
& \multicolumn{2}{X|}{Tỷ trọng đóng góp (\% of count): Quantitative (Q)} \\ \hline
\textbf{How}
& \textbf{Encode} &
\textbf{Mark:} Area (vùng góc quạt) \newline
\textbf{Channel:} \newline
\hspace*{1em} Angle: Độ lớn của tỷ trọng phần trăm (0\%--100\%) \newline
\hspace*{1em} Color: Phân biệt 2 nhóm chủ nhà \\ \cline{2-3}
& \textbf{Manipulate} &
Hover + Tooltip: Rê chuột vào từng phần quạt để xem chi tiết tên nhóm và tỷ lệ \% chính xác. \\ \cline{2-3}
& \textbf{Reduce} &
Không áp dụng. \\ \hline
\textbf{Why} & \multicolumn{2}{X|}{
\textbf{produce (derive) $\rightarrow$ explore $\rightarrow$ summarize} \newline
Derive: phân loại chủ nhà dựa trên số lượng listing sở hữu (1 phòng = Cá nhân, $>$ 1 phòng = Chuyên nghiệp). Explore và Summarize: xem nhóm nào đang chiếm ưu thế kiểm soát nguồn cung trên thị trường.
} \\ \hline
\textbf{Scale} & \multicolumn{2}{X|}{
Main key: 2 (Cá nhân, Chuyên nghiệp). \newline
Value range: 0\%--100\%.
} \\ \hline
\end{tabularx}
\end{table}

\begin{figure}[H]
    \centering
    \safeincludegraphics[width=0.9\linewidth]{charts/domain_task9_chart1.png}
    \caption{Hình 9.1. Biểu đồ tỉ trọng cơ cấu loại chủ nhà}
    \label{fig:domain-task9-chart1}
\end{figure}

\noindent\small\textit{Xem biểu đồ tương tác: }\href{https://public.tableau.com/app/profile/duy.le4605/viz/Group10_17788537028430/Task9_1-Proportionofdifferenttypesofhomeowners?publish=yes}{Tableau Public}\normalsize

\subsubsection{Đánh giá biểu đồ}

\textbf{1. Nguyên lý biểu đạt (Expressiveness):}

\begin{itemize}
    \item Thuộc tính: Loại chủ nhà (C), Tỷ trọng (Q).
    \item Channel:
    \begin{itemize}
        \item Angle: Thể hiện tỉ lệ phần trăm $\rightarrow$ dùng cho thuộc tính Q mang tính part-to-whole.
        \item Color (Hue): Phân loại chủ nhà $\rightarrow$ dùng cho thuộc tính Categorical (C).
    \end{itemize}
    \item Đánh giá mức độ quan trọng:
    \begin{itemize}
        \item Ưu tiên 1 -- Góc (Angle): Trong biểu đồ tròn, thuộc tính định lượng (Tỷ trọng \%) được ưu tiên hàng đầu để thấy được độ lớn thị phần. Theo Mackinlay, kênh Angle nằm ở mức khá cho dữ liệu (Q) và là chuẩn mực cho cấu trúc part-to-whole (thành phần trên tổng thể).
        \item Ưu tiên 2 -- Màu sắc (Hue Color): Thuộc tính phân loại (Loại chủ nhà -- C) được gán vào kênh Hue. Theo Mackinlay, Hue là kênh mạnh nhất để biểu diễn dữ liệu Categorical, giúp mắt người phân định ngay lập tức 2 nhóm mà không cần đọc text.
    \end{itemize}
\end{itemize}

\textbf{2. Nguyên lý hiệu quả (Effectiveness):}

\begin{itemize}
    \item Độ chính xác (Accuracy): Kênh Area/Angle có độ chính xác không cao. Việc ước lượng định lượng thông qua Góc/Diện tích có $n \approx 0.7$, độ lỗi Log\_error khá lớn. Mắt người khó so sánh chuẩn xác nếu 2 phần quạt có kích thường gần bằng nhau. Tuy nhiên, ở đây mức chênh lệch khá rõ (62.69\% và 37.31\%), kết hợp thao tác Manipulate (Hover xem Tooltip) đã bù đắp hoàn toàn nhược điểm này.
    \item Khả năng phân biệt (Discriminability): Biểu đồ sử dụng 2 màu tương phản (Xanh/Cam). Số lượng phần loại $< 7$ nên nhận thức màu hoàn toàn không bị ảnh hưởng, rất dễ phân biệt bằng mắt thường.
    \item Khả năng tách biệt (Separability): Kênh góc (Angle) và kênh màu sắc (Color) tách biệt hoàn toàn, không gây nhiễu cho nhau trong quá trình cảm nhận.
\end{itemize}

\subsubsection{Phân tích biểu đồ}

\begin{itemize}
    \item Thị trường Airbnb NYC mang tính thương mại hóa cao khi các chủ nhà ``Chuyên nghiệp'' (sở hữu $>$1 phòng) áp đảo nguồn cung với 62.69\% thị phần.
    \item Chủ nhà ``Cá nhân'' (Kinh doanh phòng lẻ) chỉ chiếm phần nhỏ với 37.31\%.
\end{itemize}
